\documentclass[10pt,a4paper]{amsart}
\usepackage[margin=19mm]{geometry}
\usepackage{amsmath,amssymb,mathtools}
\usepackage[scr=rsfs,cal=euler]{mathalfa}
\usepackage{xfrac}
\usepackage[unicode,hidelinks,bookmarks=false]{hyperref}
\newcommand{\ie}{i.e.\ }
\newcommand{\cf}{cf.\ }
\newcommand{\resp}{respectively\ }
\newcommand{\Subsection}{\subsection}
\providecommand{\nobreakdash}{\nobreak\hbox{-}}
\setlength{\emergencystretch}{2em}
\multlinegap0pt
\usepackage{bm}
\usepackage{bbm}
\usepackage{dsfont}
\usepackage{xspace}
\usepackage{cancel}
\usepackage{mathtools}
\usepackage{amsthm}
\makeatletter
\@ifundefined{lemma}{\newtheorem{lemma}{Lemma}}{}
\@ifundefined{theorem}{\newtheorem{theorem}{Theorem}}{}
\makeatother
\DeclareMathOperator{\R}{\mathbb{R}}
\DeclareMathOperator{\rank}{\operatorname{rank}}
\DeclareMathOperator{\spanop}{span}
\title[Tracking Through Decoupling Singularities: A Singularity-Rob]{Tracking Through Decoupling Singularities: A Singularity-Robust Homotopy-Continuation Extension of Feedback Linearization}
\author{Alex Borisevich}
\date{Source: arXiv:2607.10436v3 (CC BY 4.0)}
\begin{document}
\maketitle

\noindent\emph{Source and licence.} Tracking Through Decoupling Singularities: A Singularity-Robust Homotopy-Continuation Extension of Feedback Linearization. Version arXiv:2607.10436v3 (CC BY 4.0), distributed under the Creative Commons Attribution 4.0 International licence, as recorded in the arXiv licence metadata for this exact version and shown on the abstract page https://arxiv.org/abs/2607.10436v3. Full rights notice: licenses/arxiv-2607.10436-CC-BY-4.0.txt.\par\smallskip

\noindent\textit{Editorial scope.} Counted unit: the Theorem whose source label is \texttt{thm:crossing} in the article (source0.tex lines 678--694, source label \texttt{thm:crossing}), delivered together with the article's own complete original proof of it (source0.tex lines 696--712, one \texttt{proof} environment of 17 lines). No mathematics is written, repaired or re-derived: statement and proof are the article's own text line for line. Required context retained uncounted: lem:kernel (lines 282--299); flow (lines 638--640); error\_ode (lines 660--662). Every definition, display and label of those lines is retained unchanged. Omitted from this package and not redistributed: 1--281, 300--637, 641--659, 663--677, 695--695, 713--1913. Nothing inside a retained excerpt is omitted or altered: every retained body is delivered exactly as the article's lines are written, so no omission receipt of any kind accompanies this package. Citations: the retained excerpts carry no citation command at all, so no bibliography entry of the article is transcribed and no reference is redistributed. Reference adaptation: none was needed and none was made; every reference inside the delivered unit resolves inside it, so no unresolved reference or citation is produced. Printed slips kept literal: no printed word, symbol, hypothesis or number of the counted statement or of its proof was altered. Layout and class: the article's own class is \texttt{article} with packages including fontenc, babel, csquotes, amsmath, graphicx, amssymb, tikz, pgfplots, hyperref, geometry, amsthm, color, ifpdf, pstricks; the wrapper is the lane's pinned \texttt{amsart} editorial wrapper at 10pt on A4, which restates the theorem-like environments the retained text needs and adds the article's own macro definitions that the retained text needs (\texttt{\textbackslash R}, \texttt{\textbackslash rank}, \texttt{\textbackslash spanop}), with extra wrapper packages (bm, bbm, dsfont, xspace, cancel, mathtools). The delivered numbering is the unit's own. Assets: no retained span contains a figure environment, a graphics-inclusion command, a PDF-page inclusion, a table or a TikZ picture; no image file is copied into this package, the package declares no asset dependency and no image file is redistributed with it. Licence: version arXiv:2607.10436v3 of this article is distributed under the Creative Commons Attribution 4.0 International licence, as recorded in the arXiv licence metadata for this exact version and shown on the abstract page https://arxiv.org/abs/2607.10436v3. A full rights notice is delivered at licenses/arxiv-2607.10436-CC-BY-4.0.txt. Characters: every retained excerpt is pure ASCII, so the delivered engine, XeLaTeX with its default Latin Modern text font, reports no missing character.
\begin{lemma}[Rank and kernel of the augmented Jacobian]\label{lem:kernel}
Let $A = \bigl[ D f(x) \; \big| \; b \bigr] \in \R^{n\times(n+1)}$.
\begin{enumerate}
\item[(a)] If $\rank D f(x) = n$, then $\rank A = n$ and
\[
\ker A = \spanop\Bigl\{ \bigl( -[Df(x)]^{-1} b, \; 1 \bigr) \Bigr\}.
\]
The unit kernel vector oriented to a nonnegative last component has last ($\mu$-)component
\[
\rho \;=\; \bigl( 1 + \| [Df(x)]^{-1} b \|^2 \bigr)^{-1/2} \in (0,1].
\]
\item[(b)] If $\rank D f(x) = n-1$ and $x$ is a transversal singularity, then $\rank A = n$ (so the continuation tangent stays well defined) and
\[
\ker A = \ker D f(x) \times \{0\}.
\]
Consequently the last component of every kernel vector vanishes: $\rho = 0$.
\end{enumerate}
\end{lemma}

\begin{equation}\label{flow}
\dot c = -A^\dagger B + k\, e\, \hat n, \qquad e = t - \mu, \quad k > 0,
\end{equation}

\begin{equation}\label{error_ode}
\dot e \;=\; 1 - d(x,t) - k\, \rho(x,t)\, e .
\end{equation}

\begin{theorem}[Bounded crossing and recovery of tracking]\label{thm:crossing}
Under assumptions (A1)--(A3), the continuation flow \eqref{flow} satisfies:
\begin{enumerate}
\item[1.] (No blow-up.) The right-hand side of \eqref{flow} is finite for all $t$, including at $t^*$; in particular $\|\dot x(t)\| \le \|A^\dagger B\| + k|e|$ stays bounded through the singularity, and near $t^*$ the injected motion $k e\, \hat n$ is directed along $\ker D f(x)$.
\item[2.] (Contraction away from the singularity.) On any interval where $\rho \ge \rho_0$,
\[
|e(t)| \;\le\; |e(t_1)|\, e^{-k \rho_0 (t - t_1)} \;+\; \frac{\bar D}{k \rho_0}.
\]
In particular the output error obeys $\| f(x) - y^* \| = |e|\, \|b\| = O(1/k)$ in the regular regime.
\item[3.] (Bounded excursion during the crossing.) For all $t \in \overline{I_\delta}$,
\[
|e(t)| \;\le\; |e(t^*-\delta)| \;+\; 2\delta\, \bar D ,
\]
so the transient growth of the tracking error while feedback authority is lost is at most $2\delta\bar D$ in magnitude, and the output error stays bounded by $\|b\|_\infty\bigl(|e(t^*-\delta)| + 2\delta\bar D\bigr)$.
\item[4.] (Re-locking.) For $t \ge t^* + \delta$, item~2 applies from the value $e(t^*+\delta)$; hence $e(t)$ re-converges exponentially to the $O(\bar D/(k\rho_0))$ band. The flow returns to $\mu \approx t$ and reference tracking is recovered after the crossing.
\end{enumerate}
\end{theorem}

\begin{proof}
\emph{Item 1.} By Lemma~\ref{lem:kernel}(b) and (A2), $\rank A = n$ at $t^*$, so $A^\dagger = A^T(AA^T)^{-1}$ is bounded and $\hat n$ is a unit vector; thus the right-hand side of \eqref{flow} is finite. As $\rho \to 0$, Lemma~\ref{lem:kernel}(b) gives $\hat n \to (v, 0)$ with $v \in \ker Df(x)$, so the injected term $ke\,\hat n$ aligns with $\ker Df(x) \times \{0\}$. The velocity bound is immediate from \eqref{flow}.

\emph{Item 2.} On the interval, \eqref{error_ode} reads $\dot e = -k\rho\, e + (1-d)$ with $k\rho \ge k\rho_0 > 0$ and $|1-d| \le \bar D$. The comparison principle for the scalar linear equation with nonnegative time-varying rate gives
\[
|e(t)| \le |e(t_1)| e^{-\int_{t_1}^t k\rho\, d\tau} + \int_{t_1}^t e^{-\int_\tau^t k\rho}\,\bar D\, d\tau \le |e(t_1)| e^{-k\rho_0 (t-t_1)} + \frac{\bar D}{k\rho_0}.
\]
The output identity $f(x) - y^* = e\, b$ (from $H = 0$) yields the norm bound.

\emph{Item 3.} Since the orientation is fixed so that $\rho \ge 0$, the feedback term is dissipative. Using \eqref{error_ode},
\[
\frac{d}{dt}\,\tfrac{1}{2} e^2 \;=\; e\,\dot e \;=\; e(1-d) - k\rho\, e^2 \;\le\; |e|\,\bar D ,
\]
hence $\frac{d}{dt}|e| \le \bar D$ wherever $e \ne 0$. Integrating over $\overline{I_\delta}$ (length $2\delta$) gives $|e(t)| \le |e(t^*-\delta)| + 2\delta\bar D$. The output bound follows from $f(x)-y^* = e\, b$.

\emph{Item 4.} Apply item~2 with $t_1 = t^*+\delta$.
\end{proof}

\end{document}
